% The four stages of a bootstrapped gate and the two parts of the evaluation key.
\begin{figure}[!htb]
\centering
\begin{tikzpicture}[
    font=\footnotesize, >=stealth,
    stage/.style={draw, rounded corners=2pt, line width=0.7pt, fill=#1,
                  minimum height=1.05cm, text width=1.75cm, align=center,
                  inner sep=2pt},
    stage/.default=cBlue!10,
    kbox/.style={draw, densely dashed, rounded corners=2pt, fill=cGray!8,
                text width=3.3cm, align=center, font=\footnotesize,
                inner sep=3pt},
    lbl/.style={font=\scriptsize, align=center, cGray!90!black},
    flow/.style={->, line width=0.8pt}]
  \node[stage=cGray!6, text width=0.9cm] (in) {$c_1, c_2$};
  \node[stage, right=1.0cm of in]  (lin) {linear\\combination};
  \node[stage, right=1.0cm of lin] (br)  {blind\\rotation};
  \node[stage, right=1.0cm of br]  (se)  {sample\\extraction};
  \node[stage, right=1.0cm of se]  (ks)  {key\\switching};
  \node[stage=cGray!6, text width=0.9cm, right=1.0cm of ks] (out) {$c$};
  \draw[flow] (in)  -- node[lbl, above] {LWE} node[lbl, below] {$s,\,n$} (lin);
  \draw[flow] (lin) -- node[lbl, above] {LWE} node[lbl, below] {$s,\,n$} (br);
  \draw[flow] (br)  -- node[lbl, above] {TRLWE} node[lbl, below] {$s',\,N$} (se);
  \draw[flow] (se)  -- node[lbl, above] {LWE} node[lbl, below] {$s',\,N$} (ks);
  \draw[flow] (ks)  -- node[lbl, above] {LWE} node[lbl, below] {$s,\,n$} (out);
  \node[kbox, below=0.75cm of br] (bk) {bootstrapping key\\$n$ TRGSW ciphertexts of $s_i$ under $s'$};
  \node[kbox, below=0.75cm of ks] (ksk) {key-switching key\\$N t\,B_{\mathrm{ks}}$ LWE ciphertexts under $s$};
  \draw[->, densely dashed, cGray] (bk) -- (br);
  \draw[->, densely dashed, cGray] (ksk) -- (ks);
  \node[lbl, above=1pt of lin] {cheap, adds noise};
  \node[lbl, above=1pt of br] {$>95\%$ of the time};
  \node[lbl, above=1pt of se] {free, no noise};
  \node[lbl, above=1pt of ks] {changes key};
\end{tikzpicture}
\caption{The path of a bootstrapped gate. The two input ciphertexts are
combined linearly; blind rotation computes the gate on the encrypted phase and
gives a TRLWE ciphertext with fresh noise; sample extraction reads its constant
coefficient as an LWE ciphertext under the key $s'$, of size $N$; and key
switching brings it back to the key $s$, of size $n$, so that it can be used by
the next gate. The dashed boxes are the two parts of the public evaluation key
that the server has.}
\label{fig:pipeline}
\end{figure}
